\documentclass[12pt,a4paper]{article}
\usepackage[utf8]{inputenc}
\usepackage[T1]{fontenc}
\usepackage[french]{babel}
\usepackage{geometry}
\usepackage{fancyhdr}
\usepackage{amsmath}
\usepackage{listings}
\usepackage{xcolor}
\usepackage{hyperref}
\usepackage{graphicx}
\usepackage{enumitem}
\usepackage{tcolorbox}

\geometry{margin=2.5cm}
\setlength{\headheight}{14.5pt}

\pagestyle{fancy}
\fancyhf{}
\fancyhead[L]{Examen Blanc 1 - Administration Unix}
\fancyhead[R]{\thepage}
\fancyfoot[C]{}

\lstset{
    language=bash,
    basicstyle=\ttfamily\small,
    keywordstyle=\color{blue},
    commentstyle=\color{gray},
    stringstyle=\color{red},
    numbers=left,
    numberstyle=\tiny\color{gray},
    breaklines=true,
    frame=single,
    showstringspaces=false
}

\definecolor{questionbox}{RGB}{255,250,240}
\newtcolorbox{questionbox}{
    colback=questionbox,
    colframe=orange!50!black,
    title=Question,
    fonttitle=\bfseries
}

\definecolor{answerbox}{RGB}{240,255,240}
\newtcolorbox{answerbox}{
    colback=answerbox,
    colframe=green!50!black,
    title=Réponse,
    fonttitle=\bfseries
}

\title{\textbf{Examen Blanc 1 - Administration Unix/Linux}\\
\large Durée : 2 heures - Documents autorisés}
\author{}
\date{}

\begin{document}

\maketitle
\thispagestyle{fancy}

\section*{Instructions}
\begin{itemize}
    \item Durée : 2 heures
    \item Documents autorisés : Résumés de cours et TP
    \item Toutes les questions sont obligatoires
    \item Soignez la présentation et l'orthographe
    \item Pour les commandes, précisez si elles nécessitent des privilèges root
\end{itemize}

\newpage

% ============================================================
% QUESTION 1 : HIÉRARCHIE DES FICHIERS
% ============================================================

\section{Question 1 : Hiérarchie des Systèmes de Fichiers (10 points)}

\begin{questionbox}
\subsection*{Partie A : Définitions (4 points)}
\begin{enumerate}[label=\alph*)]
    \item Qu'est-ce que le FHS ? Expliquez son rôle et son importance.
    \item Quelle est la différence entre /bin et /usr/bin ? Pourquoi cette distinction existe-t-elle ?
\end{enumerate}

\subsection*{Partie B : Commandes (3 points)}
Donnez les commandes pour :
\begin{enumerate}[label=\alph*)]
    \item Lister tous les fichiers (y compris cachés) du répertoire /etc avec leurs permissions
    \item Trouver tous les fichiers de plus de 100 Mo dans /home
    \item Afficher les 20 dernières lignes du fichier /var/log/messages en temps réel
\end{enumerate}

\subsection*{Partie C : Répertoires (3 points)}
Expliquez le rôle et le contenu des répertoires suivants :
\begin{itemize}
    \item /var/log
    \item /proc
    \item /dev
\end{itemize}
\end{questionbox}

\newpage

% ============================================================
% QUESTION 2 : TYPES DE FICHIERS
% ============================================================

\section{Question 2 : Types de Fichiers Linux (8 points)}

\begin{questionbox}
\subsection*{Partie A : Inodes (4 points)}
\begin{enumerate}[label=\alph*)]
    \item Qu'est-ce qu'un inode ? Que contient-il ?
    \item Pourquoi le nom d'un fichier n'est-il pas stocké dans l'inode ?
    \item Comment vérifier le nombre d'inodes disponibles sur un système de fichiers ?
\end{enumerate}

\subsection*{Partie B : Liens (4 points)}
\begin{enumerate}[label=\alph*)]
    \item Quelle est la différence entre un lien physique (hard link) et un lien symbolique (soft link) ?
    \item Donnez la commande pour créer un lien symbolique nommé "python" pointant vers /usr/bin/python3
    \item Un lien symbolique peut-il pointer vers un fichier sur un autre système de fichiers ? Et un lien physique ?
\end{enumerate}
\end{questionbox}

\newpage

% ============================================================
% QUESTION 3 : DÉMARRAGE ET GRUB
% ============================================================

\section{Question 3 : Démarrage du Système et GRUB (12 points)}

\begin{questionbox}
\subsection*{Partie A : Processus de démarrage (6 points)}
Décrivez les étapes du processus de démarrage d'un système Linux, de l'alimentation jusqu'au système prêt. Mentionnez le rôle de chaque composant.

\subsection*{Partie B : GRUB2 (4 points)}
\begin{enumerate}[label=\alph*)]
    \item Pourquoi ne doit-on jamais éditer directement le fichier /boot/grub2/grub.cfg ?
    \item Après avoir modifié /etc/default/grub, quelle commande faut-il exécuter et pourquoi ?
    \item Qu'est-ce que l'initramfs et dans quels cas est-il nécessaire ?
\end{enumerate}

\subsection*{Partie C : Arrêt du système (2 points)}
Donnez les commandes pour :
\begin{enumerate}[label=\alph*)]
    \item Arrêter le système immédiatement
    \item Programmer un redémarrage dans 30 minutes avec un message d'avertissement
\end{enumerate}
\end{questionbox}

\newpage

% ============================================================
% QUESTION 4 : GESTION DES SERVICES
% ============================================================

\section{Question 4 : Gestion des Services (10 points)}

\begin{questionbox}
\subsection*{Partie A : Systemd (6 points)}
\begin{enumerate}[label=\alph*)]
    \item Quels sont les avantages de systemd par rapport à SysVinit ?
    \item Donnez les commandes pour :
    \begin{itemize}
        \item Démarrer le service httpd
        \item Vérifier si le service httpd est actif
        \item Activer le service httpd au démarrage
        \item Consulter les logs du service httpd en temps réel
    \end{itemize}
    \item Quelle est la différence entre \texttt{systemctl start} et \texttt{systemctl enable} ?
\end{enumerate}

\subsection*{Partie B : Targets (4 points)}
\begin{enumerate}[label=\alph*)]
    \item Qu'est-ce qu'un target systemd ? Donnez des exemples.
    \item Comment passer du mode graphique au mode multi-utilisateurs texte ?
    \item Quelle commande permet de connaître le target par défaut du système ?
\end{enumerate}
\end{questionbox}

\newpage

% ============================================================
% QUESTION 5 : GESTION DES PAQUETAGES
% ============================================================

\section{Question 5 : Gestion des Paquetages (10 points)}

\begin{questionbox}
\subsection*{Partie A : Concepts (4 points)}
\begin{enumerate}[label=\alph*)]
    \item Qu'est-ce qu'un paquetage ? Que contient-il ?
    \item Expliquez le concept de dépendances dans la gestion des paquets. Pourquoi est-ce important ?
    \item Quelle est la différence entre RPM et YUM ?
\end{enumerate}

\subsection*{Partie B : Commandes YUM (4 points)}
Donnez les commandes YUM pour :
\begin{enumerate}[label=\alph*)]
    \item Rechercher un paquet contenant le fichier /usr/bin/python3
    \item Installer le paquet httpd et toutes ses dépendances
    \item Lister tous les paquets installés
    \item Vérifier s'il y a des mises à jour disponibles
\end{enumerate}

\subsection*{Partie C : Dépôts (2 points)}
\begin{enumerate}[label=\alph*)]
    \item Où sont stockés les fichiers de configuration des dépôts YUM ?
    \item Comment lister tous les dépôts actifs ?
\end{enumerate}
\end{questionbox}

\newpage

% ============================================================
% QUESTION 6 : GESTION DES UTILISATEURS
% ============================================================

\section{Question 6 : Gestion des Utilisateurs et Groupes (12 points)}

\begin{questionbox}
\subsection*{Partie A : Fichiers de configuration (4 points)}
\begin{enumerate}[label=\alph*)]
    \item Décrivez la structure du fichier /etc/passwd. Que signifie le champ "x" dans la colonne des mots de passe ?
    \item Pourquoi les mots de passe sont-ils stockés dans /etc/shadow et non dans /etc/passwd ?
    \item Quelles sont les permissions du fichier /etc/shadow et pourquoi ?
\end{enumerate}

\subsection*{Partie B : Création et gestion (5 points)}
\begin{enumerate}[label=\alph*)]
    \item Donnez la commande complète pour créer un utilisateur "etudiant" avec :
    \begin{itemize}
        \item UID 2000
        \item Groupe primaire "users"
        \item Répertoire personnel /home/etudiant
        \item Shell /bin/bash
    \end{itemize}
    \item Comment ajouter l'utilisateur "etudiant" au groupe "wheel" ?
    \item Comment verrouiller le compte d'un utilisateur sans le supprimer ?
\end{enumerate}

\subsection*{Partie C : Permissions (3 points)}
\begin{enumerate}[label=\alph*)]
    \item Expliquez ce que signifie la permission 755 en mode octal
    \item Quelle est la différence entre setuid, setgid et sticky bit ?
    \item Donnez la commande pour donner les permissions rwxr-xr-x à un fichier
\end{enumerate}
\end{questionbox}

\newpage

% ============================================================
% QUESTION 7 : GESTION DES DISQUES
% ============================================================

\section{Question 7 : Gestion des Disques (10 points)}

\begin{questionbox}
\subsection*{Partie A : Partitionnement (4 points)}
\begin{enumerate}[label=\alph*)]
    \item Quelle est la différence entre MBR et GPT ? Quelles sont les limitations du MBR ?
    \item Donnez les commandes pour :
    \begin{itemize}
        \item Lister tous les disques et partitions du système
        \item Obtenir l'UUID d'une partition
    \end{itemize}
\end{enumerate}

\subsection*{Partie B : Systèmes de fichiers (3 points)}
\begin{enumerate}[label=\alph*)]
    \item Qu'est-ce que la journalisation dans un système de fichiers ? Quels sont ses avantages ?
    \item Donnez la commande pour formater la partition /dev/sdb1 en ext4
\end{enumerate}

\subsection*{Partie C : Montage (3 points)}
\begin{enumerate}[label=\alph*)]
    \item Expliquez le rôle du fichier /etc/fstab
    \item Donnez un exemple de ligne /etc/fstab pour monter /dev/sdb1 sur /mnt/data en ext4 avec montage automatique au démarrage
    \item Quelle commande permet de monter tous les systèmes de fichiers définis dans /etc/fstab ?
\end{enumerate}
\end{questionbox}

\newpage

% ============================================================
% QUESTION 8 : GESTION DU SWAP
% ============================================================

\section{Question 8 : Gestion de l'Espace Swap (8 points)}

\begin{questionbox}
\subsection*{Partie A : Concepts (4 points)}
\begin{enumerate}[label=\alph*)]
    \item Qu'est-ce que le swap ? Quel est son rôle ?
    \item Expliquez le mécanisme de pagination et le rôle du swap
    \item Quelle est la différence entre une partition swap et un fichier swap ?
\end{enumerate}

\subsection*{Partie B : Configuration (4 points)}
Décrivez les étapes complètes pour créer et activer un fichier swap de 2 Go nommé /swapfile, et le rendre permanent (actif au redémarrage).
\end{questionbox}

\newpage

\section*{Fin de l'examen}

\textbf{Bon courage !}

\end{document}

